\chapter{Results and Evaluation}
\label{ch:results}

\section{Methodological premise}
\label{sec:results-premise}

Neither repository publishes a quantitative accuracy benchmark. \texttt{dometria-bulk-insert} reports no measured extraction accuracy against a labeled test set, and ColdRAG reports no Recall, NDCG, hit-rate, or other recommendation-quality metric (\S\ref{sec:cf-content-hybrid}). This chapter does not manufacture either. It reports three kinds of evidence instead: architectural validation, that each system was built to its specified scope and runs; observed behaviour, in the form of the illustrative figures each system's own documentation records; and the evaluation instruments each system was built with, together with what they measure. It is explicit throughout about the difference between an instrument that exists and a result obtained with it. \S\ref{sec:results-limitations} returns to this premise directly, since the absence itself is a finding, not a gap in reporting.

\section{dometria-bulk-insert}
\label{sec:results-bulk-insert}

\subsection{Evaluation tooling}
\label{sec:results-bulk-insert-tooling}

Two scripts instrument the extraction pipeline of \S\ref{sec:bulk-insert-block3}, and they measure different things. \texttt{scripts/eval\_extraction.py} evaluates output at the level of the individual field, and is where the out-of-vocabulary-rate signal of \S\ref{sec:bulk-insert-normalization} surfaces: the proportion of populated fields whose value the closed vocabulary could not represent. A rising rate flags a vocabulary that has drifted from the documents it is meant to cover — a failure of the offline metadata compilation of \S\ref{sec:bulk-insert-metadata-compile}, not of the extraction call that encountered it.

\texttt{scripts/measure\_variance.py} instead runs the same document through the pipeline $N$ times and compares the results, separating variance into two levels that the pipeline can fail at independently. \emph{Schema-level agreement} asks whether Block~2 (\S\ref{sec:bulk-insert-block2}) selects the same contract on every run — a flip here changes which field set is even extracted, and is the more disruptive failure. \emph{Field-level value drift} asks, given a fixed contract, how much an individual field's extracted value varies across runs. The distinction matters because the two are not equally consequential: a schema flip invalidates the entire record, while value drift on one field leaves the rest intact. Separating them is what makes the non-determinism of \S\ref{sec:llm-limitations} — repeated calls to the same LLM are not guaranteed to return identical output even at zero temperature \citep{he2025defeating} — an observable, localized property of this specific pipeline rather than an unmeasured background risk. The pipeline does pin its sampler: every LLM call is issued at temperature 0, with an optional fixed seed. What \texttt{measure\_variance.py} reports is therefore the residue that pinning the sampler does not remove, which is precisely the quantity a self-consistency mechanism (\S\ref{sec:bulk-insert-block3}) is there to absorb.

Both scripts operate on runs the pipeline actually produces; neither presupposes or requires ground-truth labels. This is consistent with the absence noted in \S\ref{sec:results-premise}: the tooling measures self-consistency and vocabulary coverage, not accuracy against a reference.

\subsection{Quality gates}
\label{sec:results-bulk-insert-quality-gates}

Independently of the runtime evaluation tooling, every commit is checked by a fixed gate run through \texttt{pre-commit}: \texttt{ruff} for linting, \texttt{mypy --strict} for static type checking, and \texttt{pytest} for the test suite. The suite mocks the Instructor client for its unit tests, so most of it runs without network access or an API key; a smaller set of integration tests is gated behind a \texttt{-m integration} marker and requires \texttt{OPENROUTER\_API\_KEY} to exercise the pipeline against a live model. This gate is a code-correctness check — it establishes that the pipeline runs and that its types are internally consistent — and is a different kind of evidence from the field- and schema-level instrumentation of \S\ref{sec:results-bulk-insert-tooling}, which characterizes what the pipeline outputs rather than whether it runs.

The project roadmap records phases 0 through 9 as complete, none left open, including the Phase~9 hardening of \S\ref{sec:bulk-insert-hardening}. This is architectural validation in the narrow sense of \S\ref{sec:results-premise}: the specified scope was built and is exercised on every commit. It carries no claim about extraction quality.

\subsection{Illustrative response values}
\label{sec:results-bulk-insert-illustrative}

\texttt{confidence} and \texttt{completeness} belong to the response contract of \S\ref{sec:bulk-insert-api}, not to an offline report: every record carries its own self-assessment, alongside the \texttt{flags} list recording agreement, grounding and normalization notes and the \texttt{review\_queue} of unfilled mandatory fields. A low score with a non-empty review queue routes the record to a human reviewer; it is not an error condition. Per-record quality assessment is thus a runtime output of the service rather than something measured about it afterwards — which is a different thing from an aggregate accuracy figure, and does not substitute for one.

The service's own mission documentation gives \texttt{confidence: 0.82} and \texttt{completeness: 0.75} as example values in a sample response. These are illustrative: they demonstrate the shape and range of the response schema described in \S\ref{sec:bulk-insert-api}, not an aggregate measurement obtained by running the pipeline over a test corpus. No corpus-level average of either score is reported anywhere in the specification, and this thesis does not construct one.

\section{ColdRAG}
\label{sec:results-coldrag}

\subsection{Performance instrumentation}
\label{sec:results-coldrag-tracker}

The \texttt{PerformanceTracker} introduced in \S\ref{sec:coldrag-stack} (\texttt{metrics.py}) instruments a request end to end: wall-clock latency, an estimated token count from a length-based heuristic — characters divided by four, not the model's own subword tokenizer (\S\ref{sec:transformer}) — so the figure approximates the unit in which cost is billed rather than reproducing it, an estimated dollar cost derived from the configurable per-million-token input/output pricing already given in \S\ref{sec:coldrag-stack} (\$0.30 in / \$2.50 out by default), and throughput in tokens per second. A per-item latency is tracked alongside the aggregate figures where a request processes more than one item. At the end of a run, \texttt{print\_report()} emits a human-readable summary and returns a structured record — \texttt{\{elapsed\_time, total\_tokens\_estimated, cost\_dollars, tps\}} — so the same measurement is available both for inspection during development and for programmatic collection.

This instrumentation measures cost and latency, not output quality: it establishes what a request costs to run, not whether the candidates it returns are the right ones. It is the infrastructural counterpart to the accuracy figures \S\ref{sec:results-premise} notes are absent, and the two are not substitutes for each other.

\subsection{Test suite}
\label{sec:results-coldrag-tests}

ColdRAG's Pytest suite (\texttt{tests/}, run via \texttt{uv run pytest tests/ -v}) covers two areas. \texttt{test\_models.py} exercises the Pydantic layer directly: the taxonomy-driven \texttt{RuleParser} of \S\ref{sec:coldrag-hitl}, the structure of the final confirmed payload, and deterministic product-identifier generation (\S\ref{sec:coldrag-ingestion}) — that hashing the same text twice yields the same identifier is a testable property, not only an architectural claim. \texttt{test\_workflow.py} exercises the conversational module's event-driven control flow (\S\ref{sec:coldrag-hitl}): event transitions between workflow steps and session TTL behavior in Redis. As with the bulk-insert quality gate of \S\ref{sec:results-bulk-insert-quality-gates}, this suite verifies that the system behaves as specified; it does not measure retrieval or matching quality.

\subsection{An illustrative end-to-end run}
\label{sec:results-coldrag-run}

One infrastructural run is reported in the system's documentation as an example of \texttt{PerformanceTracker} output: 17.50~s elapsed, 5 items processed, approximately 1{,}200 tokens, an estimated cost of \$0.00102, a throughput of 68.57 tokens per second, and 3.50~s per item. As with the confidence and completeness values of \S\ref{sec:results-bulk-insert-illustrative}, this is a single sample presented to show what the instrumentation reports, not a benchmark result: no distribution over multiple runs or query types is given, and no claim of representativeness is made for it here.

\section{Comparative synthesis and limits of the current evaluation}
\label{sec:results-synthesis}

The two systems are evaluated by parallel but non-identical means. Both are checked by a pre-commit-gated test suite establishing code correctness, and both carry runtime instrumentation reporting on their own execution — variance and vocabulary coverage for \texttt{dometria-bulk-insert}, latency and cost for ColdRAG. Neither instrument set measures the same thing across the two systems, because the systems address different problems: an extraction service is evaluated against how faithfully and consistently it transcribes a known document, while a retrieval system is evaluated against how well it ranks candidates against a query, a task for which no held-out relevance judgments exist here at all.

The two are complementary, not equivalent. Test suites answer \emph{does the system run correctly}; variance measurement answers \emph{is a given output stable across repetitions}; cost and latency tracking answers \emph{what does a request cost}. None of the three answers \emph{is the output correct}, and neither system's tooling was built to answer that question directly. Reproducibility instrumentation is, on its own terms, a substitute for correctness only up to a point: consistent output across runs (measured by \texttt{measure\_variance.py}, \S\ref{sec:results-bulk-insert-tooling}) rules out a specific failure mode — instability — without establishing that the stable value is the right one.

One design choice accounts for both absences. In both systems, verification sits inside the pipeline rather than after it: self-consistency voting over mandatory field groups and a validator rejecting any unquoted value in \texttt{dometria-bulk-insert} (\S\ref{sec:bulk-insert-block3}), evidence and logical checks before any write in ColdRAG's governance layer (\S\ref{sec:coldrag-ingestion}), and a quote-bounded LLM judge in the out-of-vocabulary resolver (\S\ref{sec:bulk-insert-normalization}) and the HITL module (\S\ref{sec:coldrag-hitl}). In-band validation rejects unsupported output at the point of generation, which reduces what an external harness would be needed to catch downstream. It also leaves no held-out record on which such a harness could be run: the evidence-grounded design of \S\ref{sec:evidence-grounded-principle} makes the pipelines defensible by construction and unbenchmarked by the same stroke.

\section{Absence of formal benchmarks as a recognized limitation}
\label{sec:results-limitations}

No component of either system is evaluated here against a held-out, labeled ground truth. For \texttt{dometria-bulk-insert} this would mean field-level accuracy against manually annotated documents; for ColdRAG it would mean a ranking metric such as Recall@$k$ or NDCG (\S\ref{sec:cf-content-hybrid}) against known correct matches, evaluated with the kind of reference-free or reference-based RAG evaluation frameworks surveyed in \S\ref{sec:rag} \citep{es2024ragas,chen2024benchmarking}. Neither exists for either system at the time of writing, and this thesis does not attempt to construct one post hoc: doing so without a properly constructed reference set and annotation protocol would produce a number without the validity a benchmark claim requires.

This is recorded here as a limitation of the evaluation presented, not of the systems themselves, and it is treated as a direction for future work in Chapter~\ref{ch:conclusions} rather than as an omission from this chapter. What the chapter does establish is narrower and load-bearing in its own right: both systems run, are exercised by an automated test suite on every change, expose evaluation and instrumentation tooling purpose-built for the failure modes their own architecture makes possible — non-determinism and vocabulary drift for one, cost and latency for the other — and neither has been evaluated, here or in its own documentation, against a formal accuracy benchmark.
